\subsection*{Semantic publishing workflow}

The need to publish taxonomic papers as structured text has long been recognized \citep{penev2010semantic, penev2011xml, miller2012cybertaxonomic}. In conventional documents, taxon names, treatments, specimen records, figures, and references are visible to readers but are not always identified in a way that computers can interpret. Structured markup gives these elements explicit roles, allowing them to be found, linked, extracted, and reused. TaxPub introduced such markup for taxonomic literature, while Pensoft's ARPHA platform allows manuscripts to be written and reviewed online and published as structured XML, HTML, and PDF \citep{penev2011xml, stoev2018arpha}. Together with persistent identifiers, accessible repositories, licences, and suitable metadata, these developments can make taxonomic publications more consistent with the FAIR principles \citep{wilkinson2016fair}.

Comparable systems have been developed outside taxonomy. Pandoc Scholar generates several publication formats from one plain-text source, whereas Manubot combines text-based writing, Git version control, and automated manuscript generation \citep{krewinkel2017formatting, himmelstein2020open}. Both follow the same principle: content should be recorded once in a structured source and reused to produce the required outputs.

Our workflow applies this principle to semantic species descriptions, which remain uncommon despite their potential for computational reuse \citep{balhoff2013semantic, deans2015phenotypes, montanaro2024computable}. Phenoscript generates computable phenotype data and readable text from the same source. GitHub records changes, and Overleaf imports the generated Markdown into the \LaTeX{} manuscript. A correction is therefore made in Phenoscript and propagated to the manuscript after regeneration and synchronization, rather than copied into several versions by hand. This reduces transcription errors, keeps the published treatment aligned with its semantic source, and preserves a shared revision history \citep{ram2013git, braga2023github}.

Wider adoption depends partly on journals accepting \LaTeX{} source files. Where a compatible template is available, separating content from presentation reduces manual reformatting and may also reduce production work and costs \citep{bar2021reproducible, krewinkel2017formatting}. Overleaf is convenient but not essential: it offers a \href{https://www.overleaf.com/user/subscription/plans}{free tier}, and some institutions provide premium access. Alternatively, authors can use a free local \LaTeX{} distribution and text editor with GitHub. What matters is that the manuscript remains text-based and version-controlled so that generated Phenoscript outputs can be included without duplication.

Phenoscript presents its own barrier because authors must learn the syntax, select appropriate ontology terms, and validate the resulting statements. Broader use will require community-curated ontologies, reusable examples, better documentation and interfaces, and training for taxonomists with little programming experience. Artificial-intelligence tools may help convert existing descriptions or prepare initial Phenoscript drafts, but their output will still require logical validation and expert review.

This study shows that semantic descriptions can be used throughout a complete taxonomic revision. Publication alone, however, does not make them readily available for comparative use. A shared repository is still needed where semantic phenotypes can be deposited, retrieved, queried, and compared across taxa and publications. Besides providing these services, such a repository could offer a common setting in which the taxonomic community develops ontologies, tests new tools, and improves standards for computable descriptions.
